% =====================================================================
% ALGORITHM: Review of Economic Research (ARER) — manuscript template
%
% Compile with XeLaTeX if your text contains Uzbek Latin (oʻ, gʻ) or
% Cyrillic; pdfLaTeX is sufficient for plain English.
%
% Do NOT put author names in this file: they go on the title page.
% =====================================================================
\documentclass[12pt,a4paper]{article}

\usepackage[margin=2.5cm]{geometry}
\usepackage{amsmath,amssymb}
\usepackage{booktabs}
\usepackage{graphicx}
\usepackage{setspace}
\usepackage[hidelinks]{hyperref}
\usepackage[round,authoryear]{natbib}
\bibliographystyle{apalike}

\onehalfspacing
\setlength{\parindent}{1.25em}

\title{\textbf{Title of the Article in English}}
\date{}                      % no date in the anonymised file
\author{}                    % intentionally empty: see the title page

\begin{document}
\maketitle

% ---------------------------------------------------------------- abstracts
\begin{abstract}
\noindent 150--250 words, one paragraph, no citations and no abbreviations.
State the purpose, the data and method, the main result and the conclusion.

\medskip
\noindent\textbf{Keywords:} keyword one; keyword two; keyword three;
keyword four; keyword five

\noindent\textbf{JEL classification:} E52; G21; O33
\end{abstract}

\section*{Annotatsiya (oʻzbekcha)}
150--250 soʻz. Tadqiqot maqsadi, maʼlumot va usul, asosiy natija hamda
xulosani bayon eting.

\noindent\textbf{Kalit soʻzlar:} birinchi; ikkinchi; uchinchi; toʻrtinchi;
beshinchi

\section*{Аннотация (на русском)}
150--250 слов. Изложите цель исследования, данные и метод, основной
результат и вывод.

\noindent\textbf{Ключевые слова:} первое; второе; третье; четвёртое; пятое

\newpage

% ------------------------------------------------------------------- body
\section{Introduction}
State the research question, why it matters, what is new here, and how the
paper is organised.

\section{Literature review}
What is already known and the specific gap this paper fills. Cite
international literature, not only regional sources \citep{acemoglu2012}.

\section{Data and methodology}
Data sources, sample, variables, model specification, identification
strategy and software, in enough detail for the work to be reproduced.

\section{Results}

\begin{table}[htbp]
  \centering
  \caption{Descriptive statistics}
  \label{tab:descriptives}
  \begin{tabular}{lrrr}
    \toprule
    Variable & Mean & Std.\ dev. & $N$ \\
    \midrule
    Example 1 & 0.00 & 0.00 & 0 \\
    Example 2 & 0.00 & 0.00 & 0 \\
    \bottomrule
  \end{tabular}

  \smallskip
  \footnotesize Source: authors' calculations. Units: \dots
\end{table}

\section{Discussion}
Interpretation, comparison with earlier findings, mechanisms and
limitations. Do not overstate causality.

\section{Conclusion}
Answer the research question and draw only the implications the evidence
supports.

% -------------------------------------------------------------- statements
\section*{Statements}
\noindent\textbf{Funding:} This research received no specific grant from any
funding agency.\\
\textbf{Conflict of interest:} The authors declare no conflict of interest.\\
\textbf{Data availability:} Data and replication code are available from the
corresponding author on reasonable request.\\
\textbf{Use of generative AI:} No generative AI tools were used in the
preparation of this article.

% ------------------------------------------------------------- references
\begin{thebibliography}{99}
\bibitem[Acemoglu and Robinson(2012)]{acemoglu2012}
Acemoglu, D. and Robinson, J.~A. (2012).
\newblock \emph{Why Nations Fail}. Crown Business, New York.

\bibitem[Autor(2015)]{autor2015}
Autor, D.~H. (2015).
\newblock Why are there still so many jobs?
\newblock \emph{Journal of Economic Perspectives}, 29(3):3--30.
\newblock \url{https://doi.org/10.1257/jep.29.3.3}

\bibitem[Karimov(2023)]{karimov2023}
Karimov, A.~A. (2023).
\newblock Raqamli iqtisodiyot va mehnat bozori [The digital economy and the
labour market].
\newblock \emph{Iqtisodiyot va Taʼlim}, 24(3):45--58. (in Uzbek)
\end{thebibliography}

\end{document}
